%%%%%%%%%%%%%%%%%%%%%%%%%%%%%%%%%%%%%%%%%%%%%%%%%%%%%%%%%%
% FILE: typography.tex
%
% PURPOSE
% ---------------------------------
% Defines the typographic configuration of the document.
%
% This includes:
%   - font family
%   - microtypographic improvements
%   - line spacing
%
%
% ARCHITECTURE ROLE
% ---------------------------------
% Part of the style layer of the template.
%
% Responsible ONLY for visual typographic behaviour.
% It must NOT define document metadata or layout logic.
%
%
% DEPENDENCIES
% ---------------------------------
% 00_config/config.tex
% 00_config/packages.tex
%
%
% NOTES FOR DEVELOPERS
% ---------------------------------
% Do not place document content here.
%
% Changes in this file affect the entire document
% typography and readability.
%
%%%%%%%%%%%%%%%%%%%%%%%%%%%%%%%%%%%%%%%%%%%%%%%%%%%%%%%%%%


% --------------------------------------------------
% FONT FAMILY
% --------------------------------------------------

% Latin Modern font family
% Professional extension of Computer Modern

\usepackage{lmodern}

% Other popular font options:
%\usepackage{times}
%\usepackage{helvet}



% --------------------------------------------------
% MICROTYPOGRAPHY
% --------------------------------------------------

% Improves character protrusion, kerning and spacing
% Produces more professional looking text blocks

\usepackage{microtype}

\microtypesetup{
    protrusion=true,
    expansion=true
}



% --------------------------------------------------
% LINE SPACING
% --------------------------------------------------

% Academic documents usually require 1.5 spacing.
% Controlled through config.tex for flexibility.

\usepackage{setspace}

%\setstretch{\documentlinespacing}

